\section{Kinematics-first proof, methods, and supporting results}
\label{app:kin}

\subsection{Proof of Proposition~\ref{prop:sheets-switches}}
\label{app:kin-proof}
\begin{proof}
  (i) On each of the finitely many stretches where ${\vect{q}}$ is continuous, the safe one-sided limits extend the lifted pair to a continuous path in the safe fibre space, so the stretch lies in one sheet. The projected stretch images are compact, and their union carries all the mass, so full support makes it $\Omega$. The visited sheets therefore cover $\Omega$, so at least $c$ labels occur, and merging consecutive equal labels leaves at least $c-1$ changes, as in Proposition~\ref{prop:topology-transitions}. (ii) The assigned regions partition the target measure, so their conditional occupation measures, weighted by $w_j$, sum to $\mu_\Omega$. Visiting the $c$ assigned regions once each through their safe lifts gives at most $c-1$ changes, and (i) gives at least $c-1$.
\end{proof}

\subsection{Three-link experimental setup}
\label{app:kin-setup}
\paragraph{Robot, obstacles, and sheets.}
A planar arm with three unrestricted revolute joints and unit link lengths deposits on the sector $\Omega=\{(r\cos\varphi,r\sin\varphi):2.35\le r\le2.7,\ |\varphi|\le45^\circ\}\subset\mathbb R^2$, in polar coordinates $(r,\varphi)$ of the tool, among three disks of radius $0.2$ centred at $(0.5,-0.5)$, $(1.2,0.4)$ and $(1.9,-0.6)$ (Figure~\ref{fig:sheets-3r}(a)). Clearance is measured for all three links against every disk. A pose counts as available, and an inverse-kinematics solution as valid, when its clearance is at least $10^{-3}$. The safe fibre space is sampled at seven equally spaced radii from $2.35$ to $2.7$, on a grid of tool angle $\varphi$ (steps of $0.15^\circ$), last-link angle (steps of $0.25^\circ$), and elbow sign. At each radius, the safe samples are labelled into connected components. Components at adjacent radii are joined when they share at least $20$ grid cells, and a sheet is a class of joined components. This gives twelve sheets: four are present at all seven radii, and eight are fragments of at most $1378$ samples at a single radius. On a $241\times13$ availability grid of $\Omega$, an exhaustive search finds that no two sheets cover $\Omega$ and three do, so $c=3$ in the sampled model. Those three are sheets~1, 2 and~3 of Figure~\ref{fig:sheets-3r}(b), which shows where each is available; the fourth main sheet is available only where sheet~1 is. The sampled availability is saved in \path{results/kinematics/availability_grid.npz}.

\paragraph{Chart, waypoints, and densities.}
Both pipelines use the same intended sector chart, which maps the unit square onto $\Omega$ linearly in $r^2$ and in $\varphi$. It is area-preserving up to a constant, so uniform source density on the square gives uniform area density. Both use $4^m$ waypoints at orders $m=3,4,5$, and two target densities: uniform, and a gradient density that relative to area rises linearly in $\varphi$, from $0.5$ at $\varphi=-45^\circ$ to $1.5$ at $\varphi=45^\circ$. The gradient is imposed analytically through the inverse of its cumulative distribution in $\varphi$. Waypoints are joined by straight chart segments, subdivided into collocation points no more than $0.004$ apart in the workspace.

\paragraph{Geometry-first baseline.}
The baseline follows the order-$m$ Hilbert path over the whole chart. At each collocation point it solves constrained inverse kinematics, warm-started from the previous solution. It accepts the solution when it lies in the current sheet, the joint step is certified collision-free, and the tool stays within $0.002$ of the chord between the two points along the linear joint step. Otherwise it retries from the nearest sampled pose of the current sheet, and keeps the retried solution if the step to it passes the same tests. If not, the arm must reset with the flow off. The baseline then picks, among the other sheets available at that point, the one that stays available furthest ahead along the path, reading the sampled availability only as far as it needs. When the retry found the current sheet's other branch and the current sheet reaches furthest, the arm moves to that branch instead; this \emph{in-sheet reset} is recorded but is not a sheet switch. The reported minimum is the fewest sheet changes over the sampled availability of every collocation point, computed by dynamic programming. The greedy switches equal this minimum in all six requests. Solving afresh at every point is a property of this baseline, not of the geometry-first order.

\paragraph{Kinematics-first preparation.}
Preparation first chooses a chain of sheets, ordered by angle, that spans the sector, with consecutive members overlapping by at least $1^\circ$ at every sampled radius. Candidate chains are tried fewest first, and a chain is rejected before any patch is built when its best handoff curve, chosen as below, keeps a clearance below $0.03$. The chosen chain is the cover found above: sheets~1, 2 and~3. Between consecutive sheets, a handoff curve, giving the handoff angle as a function of the radius, is placed inside the sampled overlap, kept $0.3^\circ$ from its edges. A bottleneck dynamic program over the sampled radii chooses the curve that maximises the smaller of the two sheets' clearances along it, moving at most $3^\circ$ between adjacent radii. The first curve runs from $-18.2^\circ$ at the inner radius to $-22.8^\circ$ at the outer radius, and the second from $-4.1^\circ$ to $6.0^\circ$. The worst clearance along them is $0.087$ and $0.054$. Each prepared patch $U_j$ is the region between its handoff curves (or the sector edge), padded by $0.3^\circ$. Its chart is linear in $r^2$ and, at each radius, linear in $\varphi$ between the patch's edges. For each patch, preparation solves inverse kinematics in the patch's sheet on a $17\times129$ grid of the chart, row by row, maximising the smallest link-to-disk distance with a weak pull towards the neighbouring solution. It stores the result as a piecewise-affine joint-pose map on the triangulated grid, and certifies every one of its $4096$ joint-space triangles collision-free; the certified lower bounds are $1.4\times10^{-4}$, $1.0\times10^{-5}$ and $5.2\times10^{-3}$. Forward kinematics of the stored map is then compared with the intended chart on a denser $33\times257$ grid. The largest deviation is $6.8\times10^{-4}$, $1.0\times10^{-3}$ and $1.4\times10^{-4}$ for the three patches, against an acceptance threshold of $0.002$. This is a sampled validation tolerance, not a certified bound on tracking error over the patch.

\paragraph{Kinematics-first requests.}
A request follows steps (3) and (4) of the construction in Section~\ref{sec:kin-sheets}. The assigned region $\Omega_j\subseteq U_j$ is the unpadded region between its handoff curves (or the sector edge), with target-mass share $w_j=\mu_\Omega(\Omega_j)$. The source is divided into $2^m$ radial rows and $2^m$ angular columns, and the columns are split among the assigned regions in proportion to their masses by largest remainder, so the waypoint count stays $4^m$. Rounding changes the shares: for uniform density, $0.259$, $0.276$ and $0.466$ become $0.25$, $0.25$ and $0.5$ at orders three and four, and $0.25$, $0.281$ and $0.469$ at order five. The source piece of each assigned region is its rows-by-columns cell grid, traversed in a generalised Hilbert order from its low-angle to its high-angle side. It is carried into $\Omega_j$ by the Knothe--Rosenblatt map of the assigned region's conditional density: the radial marginal first, then the angle given the radius. The stored map then gives joint poses without any inverse-kinematics solve. Consecutive pieces are joined through one handoff point on their shared curve, at the radius of their first cells, near the inner radius. Collocation points include every crossing of a triangle edge, so the joint path, linear between them, equals the stored map along each chart segment. The tool path is not prescribed. It is $\FK$ of this joint motion, and only approximates the intended chart: across all requests, the largest sampled deviation at collocation points is $9.1\times10^{-4}$.

The construction realises the theoretical patch lifts only approximately. Write $\Xs_j$ for the source piece assigned to $\Omega_j$ and ${\vect{\upxi}}_j\colon[0,1]\to\Xs_j$ for its traversal. Let $\Psi_{\mathrm{k},j}\colon\Xs_j\to\Qfree$ denote the composition of the Knothe--Rosenblatt map, patch chart, and stored joint-pose map. The implementation on this piece is then
\[
  {\vect{q}}_j=\Psi_{\mathrm{k},j}\circ{\vect{\upxi}}_j,\qquad {\vect{x}}_j=\FK\circ {\vect{q}}_j,
\]
an instance of \eqref{eq:kinematicsfirst}, with $\FK\circ\Psi_{\mathrm{k},j}$ approximating the intended chart rather than equal to it. Parameterising a prepared lift through workspace coordinates does not in itself determine the construction order. What makes this construction kinematics-first is that joint motion comes from the source through a prepared, certified map, and the tool path is derived from it. The prescribed tool path is not lifted point by point. The density condition of Section~\ref{sec:kinematics-density} is targeted through the analytic chart and the conditional density maps. In density form, with $\rho_{\Qs}$ the density of $\mu_{\Qs}$ with respect to volume on the submanifold, the condition $\FK_\#\mu_{\Qs}=\mu_\Omega$ reads
\[
  \rho_{\Qs}({\vect{q}})=\rho_\Omega\bigl(\FK({\vect{q}})\bigr)\,\sqrt{\det\bigl({\vect{J}}({\vect{q}})^\top{\vect{J}}({\vect{q}})\bigr)},
\]
where ${\vect{J}}=\partial\FK/\partial{\vect{q}}$ is the Jacobian matrix of $\FK$, restricted to the tangent space of the submanifold; the square root is the factor by which $\FK$ scales volume there. A small position error does not establish this condition, because a small deviation in ${\vect{x}}$ does not bound the deviation of ${\vect{J}}$. The connectors to and from the handoff points add tool travel outside the waypoint budget: $0.19$ to $1.23$, or $0.3\%$ to $11\%$ of the tool path.

\paragraph{Certification.}
For both pipelines, every continuous stretch of the joint path is certified collision-free over adaptively subdivided pieces, using a Lipschitz lower bound on clearance. The certificate establishes positive clearance; it does not establish the $10^{-3}$ margin. A \emph{reset} is a discrete posture change at a fixed tool point: a sheet switch, or an in-sheet reset of the geometry-first baseline. It is counted, but the physical transition between postures is neither planned nor certified.

\paragraph{Measurements.}
The \emph{reset count} is the number of observed resets. The \emph{coverage error} is the largest distance from a $33\times129$ target grid to the depositing path, sampled along each continuous stretch. The \emph{density error} is the total-variation distance between target shares and the shares of the original waypoints, mapped through $\FK$, on $4\times8$ area bins. It measures where waypoints land, not continuous deposition over time: connectors, tool speed, and material flow are not in the score. Timing uses wall-clock medians of five repeats on one machine with an unoptimised Python implementation, after one untimed warm-up request. A shared fibre analysis ($15.8$~s) serves both pipelines and is excluded from the planning times. \emph{Preparation} is the method-specific one-time cost ($21.0$~s: $5.2$~s for the handoff curves and $4.6$ to $6.3$~s per patch), and \emph{planning} is the time to plan and certify one request. Geometry-first planning includes its availability lookahead but not the full availability scan that gives the reported minimum. Requests change the Hilbert order or the target density while the robot, obstacles, and target stay fixed. Kinematics-first retains its prepared maps across requests. The geometry-first baseline retains only the shared fibre analysis: as implemented, its inverse-kinematics solves and certificates belong to each request's tool points.

\paragraph{Clearance.}
Certified clearance bounds are positive but conservative worst-piece values, far below the sampled minima: $0.031$ to $0.033$ for kinematics-first, and of order $10^{-7}$ to $10^{-6}$ for geometry-first. The sampled minimum clearance of geometry-first is $0.00098$ to $0.00100$, at or just below the $10^{-3}$ imposed at each collocation point, so its paths do not quite keep that margin between points.
